% !Mode:: "TeX:UTF-8"

\chapter{通信压缩方法的设计与实现}

本章在第2章FSDP通信瓶颈分析的基础上，系统阐述在后向梯度通信压缩方面的研究工作，包括模块化通信压缩框架的设计、量化类压缩算法和稀疏化类压缩算法的实现原理与工程细节。

\section{模块化通信压缩框架}

\subsection{框架设计动机}

现有文献中的大多数通信压缩算法往往与特定的训练框架或实现环境深度耦合，例如假定参数服务器架构、单一All-Reduce通信模式，或依赖自定义的梯度聚合逻辑，因而难以在不做大幅改动的情况下直接迁移到PyTorch FSDP的分片梯度通信路径之中。为便于在统一平台上对比不同压缩方法，本文围绕FSDP的communication hook机制，设计并实现了一个模块化的通信压缩框架。

\subsection{框架架构}

该框架的核心思想是将"训练主逻辑"与"通信压缩策略"解耦，使得不同压缩算法可以在不修改训练循环与模型定义的前提下自由切换与组合。在实现层面，通信压缩框架包含以下相互配合的组成部分。

在FSDP中，每个通信钩子均遵循统一的抽象接口：
\begin{equation}
\text{hook}: (\text{state}, \mathbf{g}_{\text{full}}, \mathbf{g}_{\text{shard}}) \mapsto \text{None},
\label{eq:hook-interface}
\end{equation}
其中$\text{state}$为特定算法的状态对象，$\mathbf{g}_{\text{full}} \in \mathbb{R}^{n}$为当前FSDP单元的完整扁平梯度，$\mathbf{g}_{\text{shard}} \in \mathbb{R}^{n/W}$为当前rank上待写入的梯度分片缓冲区（$W$为world size）。Hook内部负责完成"压缩$\to$通信$\to$解码$\to$写回分片"的全过程。

各压缩算法通过各自的状态类（如GradQuantState、TopKState、SketchState等）维护必要的超参数和运行时统计量，包括量化级数、稀疏率、误差反馈残差张量以及是否启用误差反馈等。

\subsection{误差反馈机制}

误差反馈（Error Feedback）是保证有偏压缩算法收敛性的关键机制。在统一的残差更新公式下，误差反馈的实现如下：
\begin{equation}
\tilde{\mathbf{g}}_t = \mathbf{g}_t + \mathbf{e}_t, \quad \mathbf{c}_t = C(\tilde{\mathbf{g}}_t), \quad \mathbf{e}_{t+1} = \tilde{\mathbf{g}}_t - D(\mathbf{c}_t),
\label{eq:error-feedback}
\end{equation}
其中$C(\cdot)$与$D(\cdot)$分别为压缩与解压算子，$\mathbf{e}_t$为上一轮残差。残差既可以按完整梯度维度维护，也可以仅按本rank分片维度维护以减少附加显存开销。

框架通过集中式构建函数\texttt{build\_comm\_hook(name, ...)}完成压缩算法的注册与选择，将不同算法名称映射到各自的状态类与hook函数对。训练入口脚本通过命令行参数\texttt{--comm-hook}选择具体方法，再调用\texttt{model.register\_comm\_hook(state, hook)}完成注册。

\section{量化类压缩算法}

量化类方法在不改变梯度维度的前提下，通过减少每个梯度分量的表示比特数来压缩通信量。本文实现的量化方法主要包括INT8对称量化、QSGD随机量化、Natural Compression以及1-bit Seide等方案。

\subsection{INT8对称量化}

在线性INT8量化中，对梯度向量$\mathbf{g} \in \mathbb{R}^{n}$选取缩放因子：
\begin{equation}
s = \frac{q_{\max}}{\max|\mathbf{g}| + \epsilon},
\label{eq:int8-scale}
\end{equation}
其中$q_{\max}$为可表示的最大整数级，$\epsilon$为防止除零的微小常数。量化与解码过程为：
\begin{equation}
\hat{\mathbf{g}} = \text{clip}(\text{round}(\mathbf{g} \cdot s), -q_{\max}, q_{\max}), \quad \tilde{\mathbf{g}} = \hat{\mathbf{g}} / s.
\label{eq:int8-quantdequant}
\end{equation}

在FSDP集成中，各rank在本地计算缩放因子并通过一次\texttt{all\_reduce(max)}获得全局一致的缩放因子后，将扁平梯度量化为int8张量，经\texttt{reduce\_scatter(sum)}聚合，再按缩放因子与world size反缩放得到平均梯度分片。该方法实现简单，可获得约4倍通信压缩比。

\subsection{QSGD随机量化}

QSGD\cite{alistarh2017qsgd}将每个梯度分量随机舍入到有限个离散级，在期望意义下保持无偏。对非零梯度向量$\mathbf{g}$，选取级数$s$，有：
\begin{equation}
Q_{\text{QSGD}}(\mathbf{g}_i) = \|\mathbf{g}\|_2 \cdot \text{sign}(\mathbf{g}_i) \cdot \xi(|\mathbf{g}_i|, s),
\label{eq:qsgd}
\end{equation}
其中$\xi(|\mathbf{g}_i|, s) \in \{0, 1/s, \ldots, 1\}$的分布由$|\mathbf{g}_i|/\|\mathbf{g}\|_2$所在的量化区间决定，保证$\mathbb{E}[Q_{\text{QSGD}}(\mathbf{g}_i)] = \mathbf{g}_i$，并对量化方差给出上界。本实现将当前FSDP单元的完整扁平梯度$\mathbf{g}$视为一个整体，只计算一次全向量$\ell_2$范数$\|\mathbf{g}\|_2$，并对每个坐标应用上述随机量化算子。通信层面，把量化结果编码为紧凑的bit串，经\texttt{all\_to\_all}交换后在本地解码。

\subsection{Natural Compression}

Natural Compression（NC）\cite{horvath2022nc}采用与浮点指数结构天然契合的非均匀码本，将每个标量随机舍入到邻近的两个二次幂之一。对$\mathbf{g}_i \ne 0$，记$l = \log_2 |\mathbf{g}_i|$，则：
\begin{equation}
Q_{\text{nat}}(\mathbf{g}_i) = \begin{cases}
\text{sign}(\mathbf{g}_i) \cdot 2^{\lfloor l \rfloor}, & \text{以概率~} p(l), \\
\text{sign}(\mathbf{g}_i) \cdot 2^{\lceil l \rceil}, & \text{以概率~} 1-p(l),
\end{cases}
\label{eq:nc}
\end{equation}
其中$p(l)$的构造保证压缩算子无偏且二阶矩增幅受控。NC本身不依赖额外的全局缩放标量，适合多机多卡环境下的去中心化实现。量化结果打包为9比特编码（1比特符号$+$8比特指数），通过\texttt{all\_to\_all}交换后解码并求和。

此外，为进一步减轻NC在低比特情形下对收敛性的影响，本文引入范数缩放机制：在量化前首先对全局梯度向量计算$\|\mathbf{g}\|_2$，再在各rank之间通过一次\texttt{all\_reduce(sum)}汇总得到全局范数$\|\mathbf{g}\|_{2,\text{global}}$，随后使用缩放因子$\text{scale} = 1/\max(\|\mathbf{g}\|_{2,\text{global}}, \epsilon)$将梯度归一化到单位球上，在归一化梯度上应用$Q_{\text{nat}}$。解码阶段完成聚合后除以scale恢复原始尺度。

\subsection{1-bit Seide量化}

1-bit量化将每个梯度元素压缩为单个符号位$\hat{\mathbf{g}}_i = \text{sign}(\mathbf{g}_i)$，并配合标量缩放参数恢复近似梯度。Seide形式\cite{seide2014onebit}进一步在每个块上维护正负两类分量的均值$\mu^+, \mu^-$，通信阶段仅发送符号与$(\mu^+, \mu^-)$，解码时根据符号从二者中选择重构值：
\begin{equation}
\tilde{\mathbf{g}}_i = \begin{cases}
\mu^+, & \text{if~} \hat{\mathbf{g}}_i = +1, \\
\mu^-, & \text{if~} \hat{\mathbf{g}}_i = -1.
\end{cases}
\label{eq:onebit-seide}
\end{equation}
在保持1-bit通信的前提下降低重构误差。由于此类方法本质上为强有偏压缩，误差反馈几乎是必需组件。该方法可实现高达32倍的压缩比。

\section{稀疏化类压缩算法}

稀疏化类方法通过减少参与通信的梯度元素数量来压缩通信量，仅对被选中的梯度元素发送索引及其取值，其余坐标在解码端视为零。

\subsection{Top-k稀疏化}

Top-k稀疏化方法保留梯度向量中绝对值最大的$k$个坐标，其余置零\cite{stich2018sparsified}：
\begin{equation}
\text{Top}_k(\mathbf{g})_i = \begin{cases}
\mathbf{g}_i, & i \text{~属于前~} k \text{~大坐标}, \\
0, & \text{否则}.
\end{cases}
\label{eq:topk}
\end{equation}
实现时通常令$k = \lfloor \rho \cdot n \rfloor$，其中$\rho$为目标稀疏率。Top-k为有偏压缩，需配合误差反馈以保证长期收敛性。在梯度通信中，Top-k仅发送选中坐标的索引与取值，经\texttt{all\_gather}后在本地合并，在$\rho \ll 1$时可显著降低通信字节数。

\subsection{Random-k随机稀疏化}

与Top-k相比，Random-k更强调无偏性。Random-k从$\{1, \ldots, n\}$中均匀随机选取$k$个分量并保留其取值，为保持无偏性对被选中分量乘以缩放因子$n/k$：
\begin{equation}
\text{Rand}_k(\mathbf{g})_i = \begin{cases}
(n/k) \cdot \mathbf{g}_i, & i \in \mathcal{S}, \\
0, & \text{否则},
\end{cases}
\label{eq:randomk}
\end{equation}
其中$\mathcal{S}$为均匀随机选取的$k$个坐标集合。Random-k的计算复杂度较低，易于并行实现，无偏性保证理论收敛，但方差与压缩率直接相关。

\subsection{Threshold-v阈值稀疏化}

Threshold-v方法通过双阈值控制稀疏化程度\cite{dutta2020threshold}。使用正、负两个阈值$v^+ \ge 0$、$v^- \ge 0$，仅保留满足$\mathbf{g}_i \ge v^+$或$\mathbf{g}_i \le -v^-$的分量：
\begin{equation}
\text{Thresh}(\mathbf{g})_i = \begin{cases}
\mathbf{g}_i, & \mathbf{g}_i \ge v^+ \text{~或~} \mathbf{g}_i \le -v^-, \\
0, & \text{否则}.
\end{cases}
\label{eq:thresholdv}
\end{equation}
阈值可由用户显式给定，也可根据目标稀疏率从当前梯度的经验分位数自适应估计。与Top-k相比，Threshold-v避免了完整排序，计算复杂度更低。

\subsection{Sketched-SGD}

Sketched-SGD\cite{ivkin2019communication}代表了基于sketch的整体稀疏化思想。该方法通过将梯度映射到Count Sketch数据结构中来压缩维度，并利用sketch的线性性在通信侧合并各worker的sketch，再从合并后的sketch中恢复近似top-k坐标。

设Count Sketch由深度$d_s$、宽度$w_s$以及两组哈希函数$h_j: [n] \to [w_s]$、符号函数$s_j: [n] \to \{-1, 1\}$（$j = 1, \ldots, d_s$）定义，则第$j$行的sketch向量为：
\begin{equation}
S_j(\mathbf{g})[\ell] = \sum_{i: h_j(i) = \ell} s_j(i) \cdot \mathbf{g}_i, \quad \ell = 1, \ldots, w_s.
\label{eq:sketch}
\end{equation}

利用Count Sketch的线性性，有$S(\mathbf{g}^{(1)} + \mathbf{g}^{(2)}) = S(\mathbf{g}^{(1)}) + S(\mathbf{g}^{(2)})$，因此多worker情况下可以先对各自的sketch做\texttt{all\_reduce(sum)}，再在合并后的全局sketch上恢复近似top-k坐标。典型流程包括：各rank对本地梯度计算sketch表示，通过\texttt{all\_reduce}求和得到全局sketch，在全局sketch上应用HEAVYMIX策略得到重度坐标集合，最后对这些坐标执行第二轮精确值通信。理论上可将per-worker通信复杂度降到$O(k \log n)$。

\section{本章小结}

本章系统阐述了基于FSDP通信钩子机制的模块化通信压缩框架设计，以及在该框架下实现的量化类与稀疏化类共9种压缩算法。量化类算法（INT8、QSGD、NC、1-bit Seide）通过降低每个梯度分量的表示比特数实现通信压缩；稀疏化类算法（Top-k、Random-k、Threshold-v、Sketched-SGD）通过选择性传输重要梯度分量实现通信量缩减。所有算法均通过统一的hook接口与FSDP训练流程无缝集成，并支持误差反馈机制以保证收敛性。

